% MT065 / explanation — canonical editable content.
\textbf{第一问比较训练目标。} 对\(y=\pm1\)，设\(z=yf(x)\)。LR通过\(\log(1+e^{-z})\)让真实标签的概率更高，通常再加权重正则。软间隔SVM使用\(\max(0,1-z)\)加权重范数惩罚，权衡间隔与违例。LR在有限正分数处仍有小损失；hinge在\(z\ge1\)后恰为0。两者都可以搭配特征变换，不能回答成“一个只能线性、另一个非线性”。


\textbf{第二问先找到图中的异常错分区域。} 横轴远左侧\(z\ll0\)代表模型很有把握地判错。若这类点是噪声，损失曲线越陡，它要求参数改变的力量通常越大。AdaBoost的指数曲线在左侧增长最快，影响最强。


根据本题给出的图，左侧陡峭程度依次是\textbf{AdaBoost、LR、SVM}，所以按图作答为这一排序。解释要落在曲线尾部增长，而非凭算法名称背排名。


\textbf{这里必须看清LR曲线的尺度。} 原图LR在\(z=0\)附近约为1，符合\(\log_2(1+e^{-z})\)或等价缩放。换底公式\(\log_2 a=\ln a/\ln2\)说明，这相当于把自然log损失整体乘上\(1/\ln2\)，曲线高度和斜率也都乘同一个数。


这些斜率可以直接复算。导数表示曲线在当前位置的斜率；这里用到\(\frac{d}{du}\ln u=1/u\)及\(\frac{d}{du}e^u=e^u\)，复合函数还要乘内层的导数。指数损失为\(e^{-z}\)，内层\(-z\)的导数为\(-1\)，所以



\[
L'_{\mathrm{exp}}(z)=-e^{-z}.
\]


图示LR损失为\(\ln(1+e^{-z})/\ln2\)。先对log里面的整体求导，再乘\(1+e^{-z}\)的导数：



\[
L'_{\mathrm{LR,2}}(z)
=\frac1{\ln2}\frac{-e^{-z}}{1+e^{-z}}
=-\frac1{(1+e^z)\ln2}.
\]


Hinge在\(z<1\)时就是直线\(1-z\)，斜率为\(-1\)；在\(z>1\)时恒为0，斜率也为0。接点\(z=1\)有折角，不能用单个普通导数表示。这里比较的是远左侧，取斜率绝对值衡量曲线有多陡。


补充代入\(z=-4\)：指数为\(e^4\approx54.60\)；图示LR为\(1/[(1+e^{-4})\ln2]\approx1.417\)；hinge为1，符合题图排序。当\(z\)继续变得很负，\(e^z\)趋近0，图示LR的斜率绝对值趋近\(1/\ln2\approx1.443\)。若改用自然log损失，就去掉了\(1/\ln2\)这个倍数，LR在0处的高度变为\(\ln2\approx0.693\)，远负区的斜率绝对值趋近1；原图中LR比hinge更陡的严格比较不能直接沿用。


改变损失缩放还会改变它与正则项的相对强度，此题结论限于给定图形。链式求导的完整拆解见\StudyRef{Answers}{mt067}{MT067}。


\textbf{出处：} 2025A期中 Q9，原卷第6页（10分）


\textbf{答案依据：} 2025A期中Q9；同题Mock Exam Q9的清晰题图与参考答案（PDF第5页，印刷页6）。缩放解释为补充推导。 / Sources: printed exam question and matching mock question/key; scaling analysis is supplementary.
